\begin{center}
\bfseries \fontsize{24}{12} \selectfont Abstract
\end{center}
\vspace*{1\baselineskip}

\noindent Computational drug discovery has big problems because traditional \textbf{molecular docking} methods take a lot of time and use a lot of computer power. These old methods usually rely on \textbf{fixed energy scoring} and check every possible way a molecule can fit, which makes them slow and not very accurate. To fix these issues, we created \textbf{MolVista AI}, a smart system that uses \textbf{geometric deep learning} and \textbf{neural interaction models} to predict how molecules bind and how strong that binding is. Our system uses \textbf{SE(3)-equivariant graph neural networks} to predict where a drug molecule might fit into a target protein, keeping things the same even if the shape is rotated or moved. It also includes a \textbf{Neural Affinity Predictor} that looks at features like distances between atoms, angles, and contact maps to estimate how strong the binding is. \textbf{MolVista AI} is a complete system that can automatically prepare data, predict how a molecule will fit, and rate how strong the binding is, all in one process. The platform works with common file formats like \textbf{PDB}, \textbf{SDF}, and \textbf{MOL2}, making it easy to use with existing tools in the field. When we tested MolVista AI on the \textbf{PDBbind v2020} dataset, it performed very well, predicting accurate docking positions with an average \textbf{RMSD of 3.48 angstroms} and a \textbf{correlation coefficient of 0.89} with real experimental data. Additionally, it operates faster than conventional techniques, processing each molecule complex on a GPU in just \textbf{0.3 seconds}. A wide range of people can now access sophisticated docking tools thanks to \textbf{MolVista AI}'s web-based interface, which allows users to view molecules in real time, perform batch analyses, and download comprehensive results. By combining \textbf{geometric deep learning}, \textbf{neural scoring}, and \textbf{web technology}, MolVista AI sets a new standard for \textbf{AI-based molecular docking}. It helps speed up early drug discovery and improves the clarity and reliability of research on molecular interactions.

\vspace{\baselineskip}
\noindent \textbf{Keywords: } \textit{Geometric Deep Learning, SE(3)-Equivariant Graph Neural Network, Binding Affinity Prediction, Molecular Docking, Protein–Ligand Interaction, Computational Drug Discovery, Artificial Intelligence, Structural Bioinformatics}
